\clearpage\section{REGULAMENTAÇÃO DA ORGANIZAÇÃO DIDÁTICA}
\label{anexoROD}

Recorte de trechos do ROD que versam sobre matrícula, transferência e avaliação. O documento eletrônico original está disponível para acesso em \web{https://portal.ifce.edu.br/documents/3261/ROD\_atualizado.pdf}.

{\scriptsize 
\subsection*{\footnotesize  Da Matrícula: Seções I e II, Capítulo II, Título III}
\label{anexoRODmatricula}

\begin{center}
	Seção I - Da matrícula inicial
\end{center}

Art. 75. Matrícula é o ato formal pelo qual se dá a vinculação acadêmica do estudante ao IFCE após classificação em processo seletivo e convocação conforme número de vagas disponíveis, mediante apresentação dos documentos exigidos no edital.

Art. 76. Considera-se como matrícula inicial aquela realizada no período letivo de ingresso do estudante no IFCE para os cursos técnicos (integrados, concomitantes ou subsequentes) e de graduação (bacharelado, licenciatura ou tecnologia).

Art. 77. A matrícula inicial deverá ser efetivada de forma presencial pelo candidato classificado, quando maior de 18 (dezoito) anos, ou por seu representante legal, quando menor de 18 (dezoito) anos.

§ 1º Na ausência do estudante maior de 18 anos, a solicitação poderá ser realizada por seu representante legal, desde que apresente procuração com firma reconhecida.

§ 2º Na ausência do responsável legal pelo estudante menor que 18 anos, solicitação poderá ser realizada pelo representante do responsável legal, desde que apresente procuração com firma reconhecida.

§ 3º Uma vez realizada a matrícula pelo estudante, o horário da oferta dos componentes curriculares não poderá ser alterado. Em casos excepcionais, a alteração acontecerá somente mediante autorização do gestor máximo do ensino no campus.

Art. 78. Nos cursos de graduação do IFCE, é obrigatório ao estudante se matricular em todos os componentes curriculares do primeiro semestre.

Parágrafo único: Nos demais semestres o estudante deverá cumprir, no mínimo 12 (doze) créditos, salvo a condição de concludente ou em casos especiais autorizados pela coordenadoria de curso ou, na ausência desta, do gestor máximo do ensino no campus.

\begin{center}
	Seção II - Da renovação periódica da matrícula
\end{center}

Art. 79. A renovação de matrícula é um procedimento obrigatório pelo qual o estudante confirma seu interesse em manter o vínculo acadêmico com um curso do IFCE no período letivo seguinte.

§ 1º O período letivo pode se referir a um semestre letivo ou a um ano letivo, a depender da periodicidade de oferta de disciplinas do curso.

§ 2º A renovação da matrícula de um curso com periodicidade semestral deverá ser realizada a cada semestre, enquanto que para os cursos com periodicidade anual a renovação só precisará ser realizada uma vez a cada ano letivo.

Art. 80. A renovação de matrícula para os cursos técnicos e de graduação do IFCE deve ser solicitada pelo estudante de forma on-line no sistema acadêmico da instituição, de acordo com as datas previamente definidas em calendário acadêmico.

§ 1º O processo de renovação da matrícula deverá prever uma fase para solicitar a renovação e outra para ajustar a matrícula realizada pela CCA.

§ 2º O processo de renovação da matrícula deverá ser concluído até o final do período letivo que antecede o período letivo para o qual a renovação da matrícula está sendo pleiteada.

Art. 81. O estudante, que não solicitar a renovação on-line da matrícula no prazo estabelecido, deverá comparecer à CCA no prazo de 5 (cinco) dias letivos, a contar do último dia do prazo para a renovação de matrícula, a fim de regularizar sua situação acadêmica.

Parágrafo único: O estudante que não solicitar a renovação on-line da matrícula, nem comparecer fisicamente à CCA para regularizar sua situação acadêmica, deverá ser considerado desistente do curso, tendo sua situação de matrícula alterada para ABANDONO no sistema acadêmico.

\begin{center}
	Subseção II - Da renovação nos cursos de regime de créditos por disciplina
\end{center}

Art. 85. O estudante de um curso com regime de crédito por disciplina, no momento que solicitar a renovação de matrícula, deverá indicar quais componentes curriculares deseja cursar.

Parágrafo único: Os componentes curriculares a serem cursados podem ser selecionados entre aqueles:
I. obrigatórios da matriz curricular do curso;
II. optativos da matriz curricular do curso;
III. que constam em matrizes curriculares de outros cursos técnicos subsequentes ou concomitantes, desde que haja equivalência entre os componentes e que não haja choque de horário entre eles.

Art. 86. O estudante, durante a fase de ajuste de matrícula, poderá incluir ou excluir componentes curriculares para o período letivo a ser cursado.

Art. 87. Após o período de ajuste de matrículas, não deverá ser mais permitido:
I. que o estudante inclua algum componente curricular;
II. que haja alteração de horário de disciplina.

Parágrafo único: Em casos excepcionais, a alteração acontecerá somente mediante autorização do gestor máximo do ensino no campus.

Art. 88. O processo de renovação de matrícula deverá ser por componente curricular, priorizando a seguinte ordem de ocupação de vagas:
I. componentes pendentes dos estudantes finalistas;
II. componentes curriculares do semestre regular;
III. desempenho acadêmico do estudante, expresso pelo Índice de Rendimento Acadêmico (IRA).

§ 1º Entende-se por estudantes finalistas aqueles que para concluir o curso, dependem somente das disciplinas pleiteadas na renovação da matrícula.

§ 2º O Índice de Rendimento Acadêmico (IRA) é um valor quantitativo utilizado para medir o desempenho acumulado pelo estudante, nos componentes curriculares, ao longo do desenvolvimento de um curso. O cálculo do IRA é feito através de uma média ponderada das notas de cada componente curricular, levando-se em consideração a quantidade de créditos destes na matriz curricular. Este cálculo é realizado a cada fechamento de período letivo e atualizado pelo sistema acadêmico do IFCE. Para fins de cálculo, utiliza-se a seguinte fórmula:

\begin{equation*}
	IRA=\frac{(Mf_1 x Cr_1)+(Mf_2 x Cr_2)+...+(Mf_n x Cr_n)}{(Cr_1+Cr_2+...+Cr_n)}
\end{equation*}
\vspace{0,3cm}

Onde:

\begin{itemize}
	\item MF = Média final do componente curricular
	\item Cr = Créditos do componente curricular
\end{itemize}

§ 3º As notas de componentes curriculares associadas a um período letivo em curso não serão consideradas no cálculo do IRA.

§ 4º Para efeito de cálculo do IRA estarão incluídos todos os componentes curriculares cursados pelo estudante, com exceção de disciplinas com situação de trancamento, aproveitamento ou dispensa.

\subsection*{\footnotesize  Do Ingresso: Seção I, Seção II (Subseções I, II, III e IV), Seção III, IV e V}
\label{anexoRODtransferencia}

\begin{center}
	Seção II - Do ingresso de diplomados e transferidos
\end{center}

Art. 49. O IFCE poderá receber, em todos os seus cursos, estudantes oriundos de instituições devidamente credenciadas pelos órgãos normativos dos sistemas de ensino municipal, estadual e federal.

§ 1º O IFCE não receberá estudantes oriundos de cursos sequenciais.

Art. 50. O edital para ingresso de diplomados e transferidos deverá prever a seguinte ordem de prioridade de atendimento:
I. ingressantes por transferência interna;
II. ingressantes por transferência externa;
III. ingressantes diplomados.

Art. 51. Para os que pleiteiam ingresso por transferência, deverá ser considerada a seguinte ordem de prioridade no preenchimento das vagas existentes:
I. o maior número de créditos obtidos nos componentes curriculares a serem aproveitados;
II. o maior índice de rendimento acadêmico (IRA) ou índice equivalente; e
III. a maior idade.

Art. 52. No âmbito do IFCE, o ingresso de estudantes dos cursos técnicos ou de graduação, por meio de transferência, pode ser dos seguintes tipos:
I. transferência Interna
II. transferência Externa

\begin{center}
	Subseção I - Do ingresso por transferência interna
\end{center}

Art. 53. O ingresso por transferência interna é o processo de entrada de estudante em um curso de um campus do IFCE, quando este é oriundo de outro curso do mesmo campus.

Art. 54. A transferência interna só deverá ser admitida quando:
I. houver, preferencialmente, similaridade entre o curso de origem e o pleiteado no que concerne à área de conhecimento ou eixo tecnológico;
II. atender aos pré-requisitos de escolaridade e as especificidades do curso definidos em edital, mediante comprovação;
III. o curso de origem e o curso pleiteado forem do mesmo nível de ensino.

Parágrafo único – A transferência interna só poderá ser pleiteada uma vez.

\begin{center}
	Subseção II - Do ingresso por transferência externa
\end{center}

Art. 55. O ingresso por transferência externa é o processo de entrada de estudante em um curso de um campus do IFCE, quando este é oriundo de outro campus do instituto ou de outra instituição de ensino.

Art. 56. Para ter direito à matrícula, o estudante que pleiteia o ingresso por transferência deverá:
I. comprovar que foi submetido a um processo seletivo similar ao do IFCE;
II. apresentar guia de transferência ou histórico escolar com status transferido;
III. obter aprovação em teste de aptidão específica, quando o curso pretendido o exigir.

\begin{center}
	Subseção III - Do ingresso por transferência EX OFFICIO
\end{center}

Art. 57. A transferência ex officio é a forma de atendimento ao estudante egresso de outra instituição de ensino congênere, independentemente da existência de vaga, do período e de processo seletivo, por tratar-se de servidor público federal, civil ou militar, inclusive seus dependentes, e quando requerida em razão de comprovada remoção ou transferência de ofício, acarretando mudança de domicílio para o município onde se situe a instituição recebedora, ou para a localidade mais próxima desta.

§ 1º São benficiários dessa forma de ingresso o cônjuge e os dependentes do servidor até a idade de 24 anos, como caracterizado no caput deste artigo, desde que comprovado o amparo da Lei No. 9.536, de 11 de dezembro de 1997.

§ 2º Conforme estabelecido no parágrafo único da Lei No. 9.536/97, essa regra não se aplica quando o interessado na transferência se deslocar para assumir cargo efetivo em razão de concurso público, cargo comissionado ou função de confiança.

Art. 58. A solicitação de transferência ex officio deverá ser feita mediante requerimento protocolado no campus de destino e encaminhado ao gestor máximo do ensino no campus do IFCE, sendo necessários os seguintes documentos:
I. cópia do ato de transferência ex officio ou remoção, publicado no Diário Oficial da União (DOU), ou órgão oficial de divulgação ou publicação da própria corporação;
II. declaração original da autoridade maior do órgão competente, comprovando a remoção ou transferência ex officio.

\begin{center}
	Subseção IV - Do ingresso de diplomados
\end{center}

Art. 59. Entende-se por diplomados aqueles que possuem diploma de cursos de educação profissional técnica de nível médio ou diploma de cursos de graduação.

Art. 60. O requerente deverá ser diplomado no nível respectivo ou superior ao pretendido.

Art. 61. O ingresso de diplomados deverá ser concedido mediante o atendimento em pelo menos um dos seguintes critérios abaixo relacionados, desde que estes estejam definidos em edital estabelecido pelo campus:
I. maior número de créditos a serem aproveitados no curso solicitado;
II. classificação em entrevista ou prova;
III. classificação em teste de habilidades específicas, quando o curso o exigir.

Art. 62. O requerimento para ingresso de diplomado deverá ser acompanhado dos seguintes documentos, em cópia autenticada ou com a apresentação original para conferência:
I. documento oficial de identidade com foto;
II. cadastro de pessoa física (CPF);
III. cópia autenticada de diploma ou certidão de conclusão;
IV. histórico escolar;
V. programa dos componentes curriculares cursados, autenticados pela instituição de origem;
VI. outros documentos especificados em edital.

\begin{center}
	Seção III - Do ingresso por matrícula especial
\end{center}

Art. 63. Deverá ser admitida matrícula especial, ao estudante que deseje cursar componentes curriculares nos cursos técnicos e de graduação, desde que haja vaga nos componentes curriculares constantes na solicitação e que o requerente seja diplomado no nível respectivo ou superior ao pretendido.

Art. 64. O estudante com matrícula especial poderá cursar no máximo 3 (três) componentes curriculares, podendo posteriormente aproveitá-los, caso efetive uma matrícula no IFCE.

Parágrafo único: Candidatos que possuam diploma estrangeiro de curso técnico ou de graduação e se submeteram a processo de revalidação de diplomas no IFCE, poderão cursar mais de três disciplinas, na qualidade de estudante especial, desde que seja uma recomendação da comissão avaliadora da revalidação, registrada em parecer técnico.

Art. 65. A solicitação de matrícula especial deverá ser feita mediante requerimento protocolado e encaminhado à coordenadoria do curso, nos primeiros 50 (cinquenta) dias letivos do período letivo imediatamente anterior ao que deverá ser cursado, devendo ser acompanhada dos seguintes documentos:
I. cópia do diploma para quem deseja matrícula na graduação, devidamente autenticada ou acompanhada do original;
II. cópia do diploma de conclusão do curso técnico de nível médio para quem deseja matrícula em curso técnico, devidamente autenticada ou acompanhada do original;
III. cópia do histórico escolar autenticada ou acompanhada do original.

§ 1º A coordenadoria do curso pleiteado pelo interessado deverá emitir o parecer no prazo de 30 (trinta) dias.

§ 2º Caberá à Proen encaminhar o parecer técnico ao gestor máximo do ensino no campus que, por conseguinte, deverá tomar as providências de efetivação de matrícula especial desses candidatos junto à sua CCA.

Art. 66. A matrícula especial não assegura, em qualquer hipótese, vínculo como estudante regular do IFCE.

Art. 67. O estudante com matrícula especial ficará sujeito às normas disciplinares e didático- pedagógicas, inclusive submetendo-se ao sistema de avaliação do componente curricular.

Art. 68. O estudante aprovado terá direito à declaração emitida pela CCA, constando: o componente curricular cursado, a carga horária, o período, a nota, a frequência e a ementa.

Art. 69. Em nenhuma hipótese, deverá ser permitido o ingresso informal de estudante ouvinte nos cursos do IFCE, sendo, portanto, o ingresso concedido somente ao aluno com matrícula especial, mediante documentação apresentada e parecer autorizativo.

\begin{center}
	Seção IV - Do Reingresso
\end{center}

Art. 70. O IFCE concederá, em oportunidade única, o direito de reingresso a estudantes que abandonaram o curso, nas seguintes condições:
I. terem decorridos, no máximo, 5 (cinco) anos, a contar da data em que o estudante deixou de frequentar o curso;
II. existir vaga no curso;
III. apresentar em requerimento a quitação com a biblioteca (nada consta).

Art. 71. A solicitação de reingresso deverá ser feita mediante requerimento protocolado e enviado à coordenação de curso para análise e emissão de parecer.

§ 1º Em caso de deferimento da solicitação, o coordenador do curso deverá comunicar à CCA para que o estudante seja matriculado no sistema acadêmico.

§ 2º O estudante deverá receber um novo código de matrícula e ser vinculado à matriz curricular vigente do curso no qual está reingressando.

§ 3º A forma de ingresso do estudante a ser registrada no sistema acadêmico deverá ser REINGRESSO;

§ 4º Para aproveitar os componentes curriculares cursados com a matrícula anterior, o estudante deverá solicitar o aproveitamento de componentes curriculares, de acordo com os procedimentos estabelecidos no Capítulo IV - Seção I.

Art. 72. Não deverá ser permitido o reingresso de estudantes que deixaram de frequentar o curso:
I. no primeiro semestre – para cursos com periodicidade de oferta semestral de vagas;
II. no primeiro ano – para cursos com periodicidade de oferta anual de vagas.

\begin{center}
	Seção V - Da ocupação de duas vagas em cursos do mesmo nível
\end{center}   

Art. 73. No âmbito do IFCE, em nenhuma hipótese deverá ser permitida aos estudantes de cursos de graduação, a ocupação de vagas em mais de um curso do mesmo nível de ensino.

Art. 74. Ao constatar que há estudante ocupando mais de uma vaga em cursos de graduação no IFCE, ou no IFCE e em outra instituição pública, a CCA deverá comunicar ao estudante a possibilidade de optar por uma das vagas no prazo de 5 (cinco) dias úteis, contado do primeiro dia útil posterior à comunicação.

§ 1º Caso o estudante não compareça no prazo assinalado neste artigo ou não opte por uma das vagas, a instituição providenciará o cancelamento:
I. da matrícula mais antiga, na hipótese da duplicidade ocorrer em instituições diferentes;
II. da matrícula mais recente, na hipótese da duplicidade ocorrer na mesma instituição.

§ 2º Concomitantemente ao cancelamento compulsório da matrícula na forma do disposto no § 1º deste artigo, deverá ser decretada a nulidade dos créditos adquiridos no curso cuja matrícula foi cancelada.

\subsection*{\footnotesize  Da Sistemática de Avaliação: Subseção I da Seção I, Capítulo III, Título III}
\label{anexoRODavaliacao}

\begin{center}
	Seção I - Da sistemática de avaliação
\end{center}

Art. 94. Os processos, instrumentos, critérios e valores de avaliação adotados pelo professor deverão ser explicitados aos estudantes no início do período letivo, quando da apresentação do PUD, observadas as normas dispostas neste documento.

§ 1º As avaliações devem ter caráter diagnóstico, formativo, contínuo e processual, podendo constar de:
I. observação diária dos estudantes pelos professores, durante a aplicação de suas diversas atividades;
II. exercícios;
III. trabalhos individuais e/ou coletivos;
IV. fichas de observações;
V. relatórios;
VI. autoavaliação;
VII. provas escritas com ou sem consulta;
VIII. provas práticas e provas orais;
IX. seminários;
X. projetos interdisciplinares;
XI. resolução de exercícios;
XII. planejamento e execução de experimentos ou projetos;
XIII. relatórios referentes a trabalhos, experimentos ou visitas técnicas,
XIV. realização de eventos ou atividades abertas à comunidade;
XV. autoavaliação descritiva e outros instrumentos de avaliação considerando o seu caráter progressivo.

Art. 95. Ao estudante deverá ser assegurado o direito de conhecer os resultados das avaliações mediante vistas dos referidos instrumentos, apresentados pelos professores como parte do processo de ensino e aprendizagem.

§ 1º As avaliações escritas deverão ser devolvidas; e as demais, informadas ao estudante e registradas no sistema acadêmico, logo após a devida correção em um prazo máximo de até 10 (dez) dias letivos.

§ 2º A divulgação de resultados tem caráter individual, sendo vedada a sua exposição pública, salvo em casos de haver consentimento prévio do estudante.

Art. 96. O estudante que discordar do resultado obtido em qualquer avaliação da aprendizagem poderá requerer, à coordenadoria de curso, revisão no prazo de 2 (dois) dias letivos após a comunicação do resultado.

§ 1º A revisão da avaliação deverá ser feita pelo docente do componente curricular, juntamente com o coordenador do curso.

§ 2º Caso a revisão não possa ser feita pelo professor do componente curricular, o coordenador deverá designar outro docente para tal ação.

\begin{center}
	Subseção I - Avaliação nos cursos com regime de créditos por disciplina
\end{center}

Art. 97. A sistemática de avaliação dos conhecimentos construídos, nos cursos com regime de crédito por disciplina, com periodicidade semestral, se desenvolverá em duas etapas.

§ 1º Deverá ser registrada no sistema acadêmico apenas uma nota para a primeira etapa (N1) e uma nota para a segunda etapa (N2), com pesos 2 e 3, respectivamente.

§ 2º O docente deverá aplicar, no mínimo, duas avaliações em cada uma das etapas.

§ 3º O critério para composição da nota de cada etapa, a partir das notas obtidas em cada uma das avaliações, ficará a cargo do docente da disciplina, em consonância com o estabelecido no PUD.

Art. 98. O cálculo da média parcial (MP) de cada disciplina deve ser feito de acordo com a seguinte equação:

\begin{equation*}
	X_S=\frac{2X_1+3X_2}{5}\geq 7,0
\end{equation*}
\vspace{0,3cm}

Art. 99. Deverá ser considerado aprovado no componente curricular o estudante que, ao final do período letivo, tenha frequência igual ou superior a 75\% (setenta e cinco por cento) do total de horas letivas e tenha obtido média parcial (MP) igual ou superior a:
I. 6,0 (seis), para disciplinas de cursos técnicos concomitantes e subsequentes.
II. 7,0 (sete), para disciplinas de cursos de graduação.

Parágrafo único: Os estudantes aprovados com a nota da MP não precisarão realizar a avaliação final (AF) e sua média final (MF) deverá ser igual a sua média parcial (MP).

Art. 100. Deverão fazer avaliação final (AF) o estudante de curso técnico que obtiver MP inferior a 6,0 (seis) e maior ou igual a 3,0 (três), e o estudante de graduação que obtiver MP inferior a 7,0 (sete) e maior ou igual a 3,0 (três).

§ 1º A avaliação final deverá ser aplicada no mínimo 3 (três) dias letivos após o registro do resultado da MP no sistema acadêmico.

§ 2º A avaliação final poderá contemplar todo o conteúdo trabalhado no período letivo.

§ 3º A nota da avaliação final (AF) deverá ser registrada no sistema acadêmico.

§ 4º O cálculo da média final (MF) o estudante referido no caput deverá ser efetuado de acordo com a seguinte equação:

\begin{equation*}
	MF=\frac{X_S+AF}{2}\geq 5,0
\end{equation*}
\vspace{0,3cm}

§ 5º Deverá ser considerado aprovado na disciplina o estudante que, após a realização da avaliação final, obtiver média final (MF) igual ou maior que 5,0 (cinco).
}
